\makeabstract
{
Towards the Synthesis of Borangulene for Cocrystalization with TOPA
}
{
Atropa Choi (1,2), Abby A. Moravek(1), Tristan Kim(1), Ren A. Wiscons, PhD (1)
}
{
\textsuperscript{1}Department of Chemistry, Amherst College, Amherst, MA, USA\\
\textsuperscript{2}Department of Physics, Amherst College, Amherst, MA, USA
}
{
We attempted to synthesize borangulene, a planar, boron centered analog of the bowl-shaped, nitrogen centered molecule TOPA. TOPA can interconvert between two stable forms—bowl up and bowl down—and may be used for binary data storage on the molecular scale. However, further computational analysis revealed that the flat form also corresponds to a potential well, i.e., the molecule has three stable forms. Polymerizing TOPA with borangulene through a Lewis pair bond between the nitrogen and boron centers may force the bowl up and bowl down geometries to be much more favorable, thereby shifting the potential energy curve of TOPA into a double well. 
}
{
Referencing the first total synthesis of borangulene in literature(1), we attempted a modified route. The first step combines 2-bromoresorcinol and 2-bromo-1-iodo-3-methoxybenzene in a 1:3 molar ratio in an Ullmann-type reaction. After the reagents were added to a dry 3-neck round bottom flask under nitrogen gas, the mixture was stirred for 20 hours at 110°C. The mixture was passed through a celite plug using chloroform, washed with saturated NH4Cl and brine, then separated via column chromatography with 95:5 hexanes:ethyl acetate as the eluent. A sample was taken after the liquid-liquid extractions and analyzed using gas chromatography-mass spectrometry (GCMS), along with the two starting materials. 
}
{
The Ullmann-type reaction seemed to have a poor yield, and the product was difficult to separate via column chromatography. After several unsuccessful attempts, an alternative route which takes several steps from the synthesis of TOPA was explored. The GCMS chromatogram for the reaction mixture showed an intense peak matching 2-bromo-1-iodo-3-methoxybenzene, and no significant peak matching 2-bromoresorcinol. The chromatogram for 2-bromoresorcinol showed several peaks. 
}
{
Chromatograms revealed that the limiting reagent (2-bromoresorcinol) was consumed, and that the 2-bromoresorcinol sample was quite impure, the intermediate tri-brominated resorcinol being the major species. However, the synthesis of borangulene with an Ullmann-type reaction may still be viable if 2-bromoresorcinol can be purified further. The alternative route which follows the synthesis of TOPA in its first steps should still be explored. 

(1) Yuichi Kitamoto; Suzuki, T.; Miyata, Y.; Kita, H.; Funaki, K.; Oi, S. The First Synthesis and X-Ray Crystallographic Analysis of an Oxygen-Bridged Planarized Triphenylborane. Chemical Communications 2016, 52 (44), 7098–7101. https://doi.org/10.1039/c6cc02440h.
}

\newpage
